\chapter{Worked example: Diffie--Hellman key agreement}

Following Rosulek, \emph{The Joy of Cryptography} (Diffie--Hellman). Alice sends
$g^{a}$, Bob sends $g^{b}$, and both compute the shared key $g^{ab}$. Security means
that an eavesdropper seeing the transcript $(g^{a}, g^{b})$ cannot distinguish the
real shared key from a uniform group element --- exactly the decisional
Diffie--Hellman assumption. The example makes that identification precise as an
equality of advantages.

\begin{definition}[Key-agreement distinguishing game]
  \label{def:ka-game}
  \uses{def:spcomp, def:elgamal}
  \lean{CatCrypt.Examples.DiffieHellman.KA_Game_Real, CatCrypt.Examples.DiffieHellman.KA_Advantage}
  \leanok
  The real game hands the distinguisher the transcript $(g^{a}, g^{b})$ and the
  true shared key $g^{ab}$; the ideal game replaces the key by a uniform group
  element. The advantage is the distance between them.
\end{definition}

\begin{theorem}[Key indistinguishability equals DDH]
  \label{thm:ka-ddh}
  \uses{def:ka-game}
  \lean{CatCrypt.Examples.DiffieHellman.ka_advantage_eq_ddh}
  \leanok
  The key-agreement advantage is \emph{equal} to the DDH advantage of an explicit
  reduction adversary that repackages a DDH triple $(g^{a}, g^{b}, z)$ into the
  transcript-plus-candidate-key shape:
  \[
    \mathsf{Adv}^{\mathrm{KA}}(D) \;=\; \mathsf{Adv}^{\mathrm{DDH}}(R_D).
  \]
\end{theorem}

\begin{corollary}[Security under DDH]
  \label{cor:ka-secure}
  \uses{thm:ka-ddh}
  \lean{CatCrypt.Examples.DiffieHellman.ka_secure_under_ddh}
  \leanok
  If DDH is $\varepsilon$-hard in the group, Diffie--Hellman key agreement has
  advantage at most $\varepsilon$.
\end{corollary}
